\chapter{Uzwarcenia Aleksandrowa i Čecha--Stone'a}
\label{ch:topog-uzwarcenia-kanoniczne}

\noindent\textbf{Cel rozdziału.}
W rozdziale~16 ustaliliśmy, czym jest uzwarcenie
Hausdorffa, kiedy ono istnieje i jak porównywać
uzwarcenia przez odwzorowania zgodne z zanurzeniami.
Zbudujemy teraz dwa kanoniczne uzwarcenia,
które rozwiązują różne zagadnienia:
uzwarcenie Aleksandrowa dodaje jeden punkt
do niezwartej przestrzeni lokalnie zwartej,
a uzwarcenie Čecha--Stone'a umożliwia
przedłużanie wszystkich odwzorowań
ciągłych do zwartych przestrzeni Hausdorffa.
Dla każdego z nich podamy dokładne założenia
i udowodnimy odpowiednią własność uniwersalną.

\section{Uzwarcenie Aleksandrowa}
\label{sec:topog-uzwarcenia-aleksandrow}

\begin{definicja}[Konstrukcja Aleksandrowa]
\label{def:topog-uzwarcenia-aleksandrow}
Niech $X$ będzie niezwartą przestrzenią
Hausdorffa, a $\infty$ nowym punktem,
nienależącym do $X$. Na zbiorze
$X^\ast=X\cup\{\infty\}$ rozważmy topologię,
której zbiory otwarte są następujące:
\begin{enumerate}[label=(\roman*)]
\item wszystkie zbiory otwarte w $X$;
\item zbiory $X^\ast\setminus C$,
      gdzie $C\subseteq X$ jest zwarty.
\end{enumerate}
Parę $(X^\ast,\iota)$, gdzie
$\iota\colon X\hookrightarrow X^\ast$
jest włączeniem, nazywamy
\emph{konstrukcją Aleksandrowa}.
\end{definicja}

\begin{lemat}[Topologia i zwartość konstrukcji]
\label{lem:topog-uzwarcenia-aleksandrow-topologia}
Niech $X$ będzie niezwartą przestrzenią
Hausdorffa. Powyższe rodziny tworzą
topologię na $X^\ast$,
względem której $X^\ast$ jest zwarte,
$X$ jest otwartą podprzestrzenią
z pierwotną topologią, a $X$ jest
gęste w $X^\ast$.
\end{lemat}
\begin{proof}
Zbiory $\varnothing$ i $X^\ast$ są otwarte:
pierwszy należy do (i), drugi otrzymujemy
w (ii) dla $C=\varnothing$.
Przecięcie dwóch zbiorów typu (ii)
ma postać $X^\ast\setminus(C_1\cup C_2)$,
gdzie $C_1\cup C_2$ jest zwarty.
Jeżeli $U\subseteq X$ jest otwarty,
a $C\subseteq X$ zwarty, to
$U\cap(X^\ast\setminus C)=U\setminus C$
jest otwarty w $X$, ponieważ zwarty
podzbiór przestrzeni Hausdorffa
jest domknięty.
Dowolna suma zbiorów typu (ii),
jeśli rodzina nie jest pusta,
jest dopełnieniem przecięcia
zbiorów zwartych w $X$;
przecięcie to jest zwarte,
gdyż stanowi domknięty podzbiór
jednego ustalonego zbioru zwartego.
Suma zbiorów typu (i) jest otwarta w $X$.
Suma zbioru typu (ii) z dowolnym otwartym
$U\subseteq X$ jest dopełnieniem
$C\setminus U$, które jest zwarte jako
domknięty podzbiór zwartego $C$.
Otrzymaliśmy aksjomaty topologii.

Zbiory otwarte w $X^\ast$, które
nie zawierają $\infty$, są dokładnie
zbiorami otwartymi w $X$.
Stąd $\iota$ jest zanurzeniem,
a $X=X^\ast\setminus\{\infty\}$
jest otwarty w $X^\ast$.

Niech $\mathcal V$ będzie dowolnym
pokryciem otwartym $X^\ast$.
Wybierzmy $V_\infty\in\mathcal V$
zawierający $\infty$.
Istnieje zwarty $C\subseteq X$ taki, że
$V_\infty=X^\ast\setminus C$.
Zbiory pozostałego pokrycia
indukują pokrycie otwarte $C$;
zwartość $C$ daje skończone
podpokrycie. Wraz z $V_\infty$
pokrywa ono $X^\ast$.

Gdyby jakiś otwarty zbiór
zawierający $\infty$ nie przecinał $X$,
miałby postać $\{\infty\}=X^\ast\setminus X$,
co oznaczałoby zwartość $X$.
Przeczy to założeniu. Dlatego
$\infty\in\overline X^{\,X^\ast}$
i $X$ jest gęste.
\end{proof}

\begin{twierdzenie}[Kryterium uzwarcenia Aleksandrowa]
\label{thm:topog-uzwarcenia-aleksandrow-kryterium}
Niech $X$ będzie niezwartą przestrzenią
Hausdorffa. Konstrukcja Aleksandrowa
$(X^\ast,\iota)$ jest uzwarceniem
Hausdorffa przestrzeni $X$
wtedy i tylko wtedy, gdy $X$
jest lokalnie zwarta.
\end{twierdzenie}
\begin{proof}
Załóżmy lokalną zwartość $X$.
Z poprzedniego lematu mamy
zwartość $X^\ast$ i gęstość $X$.
Dwa różne punkty $x,y\in X$
mają rozłączne otoczenia otwarte w $X$,
a więc również w $X^\ast$.
Niech $x\in X$.
Z lematu o zwartych domknięciach
z rozdziału~14 otrzymujemy otwarty
$U\ni x$ o zwartym domknięciu
$C=\overline U^{\,X}$.
Zbiory $U$ i $X^\ast\setminus C$
są rozłącznymi otoczeniami
odpowiednio $x$ i $\infty$.
Stąd $X^\ast$ jest Hausdorffa.

Odwrotnie, jeśli $X^\ast$ jest
Hausdorffa, to jako przestrzeń zwarta
Hausdorffa jest lokalnie zwarta.
Podprzestrzeń otwarta $X$
jest lokalnie zwarta,
co udowodniono w rozdziale~14.
\end{proof}

\begin{uwaga}[Dlaczego wymagamy niezwartości]
Jeżeli $X$ jest zwarta, to przy
tej samej definicji punkt $\infty$
jest izolowany, ponieważ
$\{\infty\}=X^\ast\setminus X$
jest otwarty. Wówczas $X$
nie jest gęsty w $X^\ast$,
więc konstrukcja nie jest
uzwarceniem w sensie rozdziału~16.
Przestrzeń zwarta Hausdorffa
ma już uzwarcenie $(X,\mathrm{id}_X)$.
\end{uwaga}

\section{Własność uniwersalna i minimalność}
\label{sec:topog-uzwarcenia-aleksandrow-uniwersalne}

\begin{definicja}[Odwzorowanie właściwe]
\label{def:topog-uzwarcenia-wlasciwe}
Niech $X,Y$ będą przestrzeniami
topologicznymi. Odwzorowanie ciągłe
$f\colon X\to Y$ nazywamy
\emph{właściwym}, jeżeli
$f^{-1}(C)$ jest zwarty w $X$
dla każdego zwartego $C\subseteq Y$.
\end{definicja}

\begin{twierdzenie}[Przedłużanie odwzorowań właściwych]
\label{thm:topog-uzwarcenia-aleksandrow-proper}
Niech $X,Y$ będą niezwartymi,
lokalnie zwartymi przestrzeniami Hausdorffa,
a $f\colon X\to Y$ odwzorowaniem ciągłym.
Istnieje ciągłe odwzorowanie
$f^\ast\colon X^\ast\to Y^\ast$
spełniające
$f^\ast|_X=f$ oraz
$f^\ast(\infty_X)=\infty_Y$
wtedy i tylko wtedy, gdy $f$
jest właściwe.
Przy tych warunkach $f^\ast$
jest wyznaczone jednoznacznie.
\end{twierdzenie}
\begin{proof}
Wymagane wartości jednoznacznie
wyznaczają odwzorowanie zbiorów.
Załóżmy, że $f$ jest właściwe.
Dla otwartego $V\subseteq Y$
przeciwobraz $(f^\ast)^{-1}(V)
=f^{-1}(V)$ jest otwarty w $X$,
więc również w $X^\ast$.
Jeśli otwarty zbiór w $Y^\ast$
zawiera $\infty_Y$, ma postać
$Y^\ast\setminus C$ dla pewnego
zwartego $C\subseteq Y$.
Jego przeciwobraz jest równy
$X^\ast\setminus f^{-1}(C)$,
otwarty z definicji konstrukcji,
bo $f^{-1}(C)$ jest zwarty.
Stąd ciągłość $f^\ast$.

Odwrotnie, niech $f^\ast$
będzie ciągłe i spełnia
podane warunki. Dla zwartego
$C\subseteq Y$ zbiór $C$ jest
domknięty w Hausdorffa $Y^\ast$.
Zatem $(f^\ast)^{-1}(C)$
jest domknięty w zwartym $X^\ast$,
a więc zwarty. Nie zawiera
$\infty_X$, ponieważ
$f^\ast(\infty_X)=\infty_Y\notin C$,
i jest równy $f^{-1}(C)$.
Jako podzbiór otwartej podprzestrzeni
$X$ jest również zwarty w $X$.
\end{proof}

\begin{twierdzenie}[Minimalność uzwarcenia Aleksandrowa]
\label{thm:topog-uzwarcenia-aleksandrow-minimalne}
Niech $X$ będzie niezwartą,
lokalnie zwartą przestrzenią Hausdorffa,
a $(K,e)$ dowolnym jej uzwarceniem
Hausdorffa. Istnieje dokładnie jedno
ciągłe odwzorowanie
$p\colon K\to X^\ast$ takie, że
$p\circ e=\iota$.
Jest ono surjektywne
i przenosi cały zbiór
$K\setminus e(X)$ na punkt $\infty$.
W notacji porządku rozdziału~16
oznacza to $(X^\ast,\iota)\preccurlyeq(K,e)$.
\end{twierdzenie}
\begin{proof}
Pokażemy najpierw, że $e(X)$
jest otwarty w $K$.
Ustalmy $x\in X$ i wybierzmy
otwarty zbiór $U\subseteq X$
taki, że $x\in U$ i
$C=\overline U^{\,X}$ jest zwarty.
Istnieje otwarty $O\subseteq K$
z $O\cap e(X)=e(U)$.
Ponieważ $e(X)$ jest gęsty w $K$,
zbiór $O\cap e(X)$ jest gęsty w $O$,
więc
$O\subseteq\overline{e(U)}^{\,K}$.
Zbiór $e(C)$ jest zwarty, a zatem
domknięty w Hausdorffa $K$;
zawiera $e(U)$, skąd
$O\subseteq e(C)\subseteq e(X)$.
Każdy $e(x)$ ma więc
otwarte w $K$ otoczenie zawarte
w $e(X)$, co dowodzi otwartości.

Zbiór $F=K\setminus e(X)$ jest
domknięty w $K$, a zatem zwarty.
Jest niepusty: w przeciwnym razie
$X$ byłaby homeomorficzna ze zwartym
$K$, wbrew założeniu.
Definiujemy
$p(e(x))=x$ dla $x\in X$
oraz $p(z)=\infty$ dla $z\in F$.
Jeśli $V\subseteq X$ jest otwarty,
to $p^{-1}(V)=e(V)$ jest otwarty w $K$,
ponieważ $e(X)$ jest otwarty.
Jeżeli $X^\ast\setminus C$ jest
otoczeniem $\infty$ dla zwartego
$C\subseteq X$, to
$p^{-1}(X^\ast\setminus C)
=K\setminus e(C)$ jest otwarty:
$e(C)$ jest zwarty, więc domknięty w $K$.
To dowodzi ciągłości $p$.
Odwzorowanie jest surjektywne,
gdyż obejmuje każdy punkt $X$
oraz $\infty$ (zbiór $F$ jest niepusty).
Jednoznaczność wynika z gęstości
$e(X)$ i Hausdorffowości $X^\ast$,
zgodnie z rozdziałem~16.
\end{proof}

\begin{wniosek}[Jednoznaczność uzwarcenia jednopunktowego]
\label{cor:topog-uzwarcenia-aleksandrow-jednoznacznosc}
Niech $X$ będzie niezwartą, lokalnie zwartą
przestrzenią Hausdorffa. Jeżeli $(K,e)$
jest jej uzwarceniem Hausdorffa i
$K\setminus e(X)$ składa się z dokładnie
jednego punktu, to $(K,e)$ jest równoważne
uzwarceniu Aleksandrowa $(X^\ast,\iota)$.
\end{wniosek}
\begin{proof}
Z poprzedniego twierdzenia istnieje
ciągłe surjektywne
$p\colon K\to X^\ast$ zgodne
z zanurzeniami. Jest ono różnowartościowe
na $e(X)$, a jedyny punkt
$K\setminus e(X)$ przechodzi na $\infty$;
zatem $p$ jest bijekcją.
Ciągła bijekcja ze zwartej
przestrzeni $K$ na przestrzeń
Hausdorffa $X^\ast$ jest
homeomorfizmem. Spełnia
$p\circ e=\iota$,
więc uzwarcenia są równoważne.
\end{proof}

\begin{przyklad}[Prosta i okrąg]
\label{ex:topog-uzwarcenia-aleksandrow-prosta}
Dla $X=\R$ funkcja
$e(t)=\exp(i(2\arctan t+\pi))$
jest homeomorfizmem $\R$
na $S^1\setminus\{1\}$.
Jej obraz jest gęsty w $S^1$.
Otrzymane uzwarcenie jest
równoważne uzwarceniu Aleksandrowa:
odwzorowanie $e$ przedłużamy,
przypisując $\infty$ wartość $1$.
Oznaczmy to przedłużenie przez
$\widetilde e\colon X^\ast\to S^1$.
Wykażemy jego ciągłość w $\infty$.
Ponieważ
$\lim_{t\to-\infty}2\arctan t=-\pi$
oraz
$\lim_{t\to+\infty}2\arctan t=\pi$,
mamy
$\lim_{t\to\pm\infty}e(t)=1$.
Dla każdego otwartego otoczenia
$V\subseteq S^1$ punktu $1$
istnieje $M>0$ takie, że
$e(t)\in V$ dla $|t|>M$.
Zatem otoczenie
$X^\ast\setminus[-M,M]$ punktu
$\infty$ jest zawarte w
$\widetilde e^{-1}(V)$.
Dowodzi to ciągłości w $\infty$;
na otwartym $X$ ciągłość
wynika z ciągłości $e$.
Przedłużenie jest ciągłą bijekcją
ze zwartej $X^\ast$ do Hausdorffa $S^1$,
więc jest homeomorfizmem.
\end{przyklad}

\begin{lemat}[Całkowita regularność przestrzeni lokalnie zwartych]
\label{lem:topog-uzwarcenia-lokalna-tych}
Każda lokalnie zwarta przestrzeń
Hausdorffa jest przestrzenią Tychonoffa.
\end{lemat}
\begin{proof}
Niech $X$ będzie lokalnie zwarta
i Hausdorffa, $F\subseteq X$ domknięty,
a $x\in X\setminus F$.
Lemat o zwartych domknięciach
z rozdziału~14 daje otwarty
$W\subseteq X$ taki, że
$x\in W$, a
$C=\overline W^{\,X}$ jest zwarty
i zawarty w $X\setminus F$.
Przestrzeń $C$ jest zwarta i Hausdorffa,
więc normalna.
Zbiory $\{x\}$ i $C\setminus W$
są domknięte i rozłączne w $C$.
Lemat Urysohna zastosowany w $C$
daje ciągłą $g\colon C\to[0,1]$
z $g(x)=1$ i $g|_{C\setminus W}=0$.
Definiujemy $h\colon X\to[0,1]$
przez $h|_C=g$ i
$h|_{X\setminus W}=0$.
Jest to definicja zgodna na części
wspólnej $C\setminus W$.
Oba zbiory $C$ i $X\setminus W$
są domknięte i pokrywają $X$,
dlatego lemat o sklejaniu funkcji
ciągłych daje ciągłość $h$.
Mamy $h(x)=1$ i $h|_F=0$,
gdyż $F\cap C=\varnothing$.
Funkcja $1-h$ jest więc równa $0$
w $x$ i $1$ na $F$, co dowodzi
całkowitej regularności.
Aksjomat $T_1$ wynika
z Hausdorffowości.
\end{proof}

\section{Uzwarcenie Čecha--Stone'a}
\label{sec:topog-uzwarcenia-cech-stone}

\begin{definicja}[Konstrukcja Čecha--Stone'a]
\label{def:topog-uzwarcenia-cech-stone}
Niech $X$ będzie przestrzenią Tychonoffa.
Oznaczmy $I=C(X,[0,1])$.
Z rozdziału~7 znamy zanurzenie
$e_X\colon X\to[0,1]^I$ dane wzorem
$e_X(x)=(f(x))_{f\in I}$.
Definiujemy
\[
\beta X=\overline{e_X(X)}^{\, [0,1]^I}.
\]
Parę $(\beta X,e_X)$ nazywamy
\emph{uzwarceniem Čecha--Stone'a}
przestrzeni $X$.
\end{definicja}

\begin{twierdzenie}[Uniwersalna własność uzwarcenia Čecha--Stone'a]
\label{thm:topog-uzwarcenia-cech-stone-uniwersalne}
Niech $X$ będzie przestrzenią Tychonoffa,
$L$ zwartą przestrzenią Hausdorffa,
a $f\colon X\to L$ odwzorowaniem ciągłym.
Istnieje dokładnie jedno
odwzorowanie ciągłe
$\widehat f\colon\beta X\to L$
takie, że $\widehat f\circ e_X=f$.
\end{twierdzenie}
\begin{proof}
Z rozdziału~7 wiemy, że
$e_L\colon L\to[0,1]^J$,
gdzie $J=C(L,[0,1])$,
jest zanurzeniem.
Ponieważ $L$ jest zwarta,
$e_L(L)$ jest zwartym, a zatem
domkniętym podzbiorem kostki
Hausdorffa $[0,1]^J$.

Każde $g\in J$ wyznacza
funkcję $g\circ f\in I$.
Definiujemy
$\Phi\colon\beta X\to[0,1]^J$
współrzędnymi
\[
\Phi(z)_g=z_{g\circ f},\qquad g\in J.
\]
Każda współrzędna jest ograniczeniem
rzutowania z kostki $[0,1]^I$,
więc $\Phi$ jest ciągła z własności
uniwersalnej topologii produktowej.
Dla $x\in X$ i $g\in J$ mamy
\[
\Phi(e_X(x))_g=(g\circ f)(x)
=e_L(f(x))_g.
\]
Stąd $\Phi(e_X(X))\subseteq e_L(L)$.
Gęstość $e_X(X)$ w $\beta X$
i domkniętość $e_L(L)$ dają
$\Phi(\beta X)\subseteq e_L(L)$.
Możemy zatem zdefiniować
$\widehat f=e_L^{-1}\circ\Phi$.
Jest ono ciągłe i spełnia
$\widehat f\circ e_X=f$.
Jeśli dwa ciągłe odwzorowania
$\beta X\to L$ są równe
na gęstym $e_X(X)$,
to są równe wszędzie,
ponieważ $L$ jest Hausdorffa.
Otrzymujemy jednoznaczność.
\end{proof}

\begin{wniosek}[Przedłużanie funkcji ograniczonych]
\label{cor:topog-uzwarcenia-cech-stone-funkcje}
Niech $X$ będzie przestrzenią Tychonoffa.
Każda ograniczona funkcja ciągła
$f\colon X\to\R$ ma dokładnie jedno
ciągłe przedłużenie
$\widehat f\colon\beta X\to\R$.
\end{wniosek}
\begin{proof}
Jeśli $f(X)\subseteq[a,b]$
dla liczb rzeczywistych $a\le b$,
to stosujemy poprzednie twierdzenie
do zwartego przedziału $[a,b]$,
a następnie składamy przedłużenie
z włączeniem $[a,b]\hookrightarrow\R$.
Jednoznaczność wynika z gęstości
$e_X(X)$ w $\beta X$.
\end{proof}

\begin{twierdzenie}[Maksymalność uzwarcenia Čecha--Stone'a]
\label{thm:topog-uzwarcenia-cech-stone-maksymalne}
Niech $X$ będzie przestrzenią Tychonoffa,
a $(K,e)$ dowolnym uzwarceniem
Hausdorffa $X$. Istnieje dokładnie jedno
ciągłe odwzorowanie
$p\colon\beta X\to K$
z $p\circ e_X=e$.
Jest ono surjektywne.
Uzwarcenie $(K,e)$ jest równoważne
$(\beta X,e_X)$ wtedy i tylko wtedy,
gdy każda ciągła funkcja
$f\colon X\to[0,1]$ ma ciągłe
przedłużenie do $K$ zgodne z $e$.
\end{twierdzenie}
\begin{proof}
Istnienie i jednoznaczność $p$
wynikają z własności uniwersalnej
dla odwzorowania $e\colon X\to K$.
Surjektywność: obraz $p(\beta X)$
jest zwarty, a więc domknięty
w Hausdorffa $K$, i zawiera
gęsty $e(X)$.

Jeśli uzwarcenia są równoważne,
to każda współrzędna
$f\colon X\to[0,1]$
przedłuża się do $\beta X$,
a homeomorfizm $\beta X\to K$
pozwala przenieść ją na $K$.

Odwrotnie, załóżmy istnienie
przedłużenia $F_f\colon K\to[0,1]$
dla każdego $f\in I=C(X,[0,1])$,
z $F_f\circ e=f$.
Niech $\pi_f\colon[0,1]^I\to[0,1]$
oznacza rzutowanie na współrzędną
indeksowaną funkcją $f\in I$.
Z własności jednoznaczności
odwzorowań ciągłych na gęstym obrazie
wynika
\[
F_f\circ p=\pi_f|_{\beta X}
\qquad(f\in I),
\]
ponieważ obie strony mają wartości
$f(x)$ w punktach $e_X(x)$.
Jeśli $p(z)=p(w)$, to
$z_f=F_f(p(z))=F_f(p(w))=w_f$
dla wszystkich $f\in I$.
Jako punkty kostki $[0,1]^I$
mamy $z=w$, więc $p$ jest iniektywne.
Ciągła bijekcja ze zwartej $\beta X$
na Hausdorffa $K$ jest homeomorfizmem.
Jest zgodna z zanurzeniami,
zatem uzwarcenia są równoważne.
\end{proof}

\section{Porównanie konstrukcji i przykłady}
\label{sec:topog-uzwarcenia-porownanie}

\begin{wniosek}[Dwa krańce porządku uzwarceń]
\label{cor:topog-uzwarcenia-krance}
Niech $X$ będzie niezwartą,
lokalnie zwartą przestrzenią Hausdorffa.
Dla każdego uzwarcenia Hausdorffa $(K,e)$
zachodzą nierówności w porządku
z rozdziału~16:
\[
(X^\ast,\iota)\preccurlyeq(K,e)
\preccurlyeq(\beta X,e_X).
\]
\end{wniosek}
\begin{proof}
Pierwsza nierówność wynika
z minimalności uzwarcenia Aleksandrowa,
druga z maksymalności uzwarcenia
Čecha--Stone'a. Lokalna zwartość i Hausdorffowość
zapewniają aksjomat Tychonoffa
na mocy
lematu~\ref{lem:topog-uzwarcenia-lokalna-tych}.
\end{proof}

\begin{przyklad}[Uzwarcenie dyskretnej przestrzeni liczb naturalnych]
\label{ex:topog-uzwarcenia-beta-n}
Niech $X=\mathbb N$ z topologią
dyskretną. Na mocy wniosku
\ref{cor:topog-uzwarcenia-aleksandrow-jednoznacznosc}
uzwarcenie Aleksandrowa
jest równoważne uzwarceniu
$\{0\}\cup\{1/n:n\in\mathbb N\}$
z rozdziału~16, ponieważ
to ostatnie dodaje do $\mathbb N$
dokładnie jeden punkt.
Każda funkcja $f\colon\mathbb N\to[0,1]$
jest ciągła, ale przedłuża się
do tego uzwarcenia dokładnie wtedy,
gdy ciąg $(f(n))$ jest zbieżny.

W szczególności funkcja
$f(n)=0$ dla parzystych $n$
oraz $f(n)=1$ dla nieparzystych $n$
nie ma takiego przedłużenia.
Ma natomiast ciągłe przedłużenie
do $\beta\mathbb N$
na mocy twierdzenia
\ref{thm:topog-uzwarcenia-cech-stone-uniwersalne}.
Z kryterium równoważności wynika,
że uzwarcenia Aleksandrowa i Čecha--Stone'a
przestrzeni $\mathbb N$ nie są równoważne.
\end{przyklad}

\begin{uwaga}[Założenia i nazewnictwo]
Uzwarcenie Aleksandrowa w postaci
uzwarcenia Hausdorffa wymagającego
dokładnie jednego nowego punktu
dotyczy niezwartych, lokalnie zwartych
przestrzeni Hausdorffa.
Uzwarcenie Čecha--Stone'a istnieje
dla dowolnych przestrzeni Tychonoffa
i jest największe w porządku uzwarceń.
Te własności są różne:
pierwsza konstrukcja kontroluje
odwzorowania właściwe, druga
wszystkie ciągłe odwzorowania
do zwartych przestrzeni Hausdorffa.
\end{uwaga}
